\documentclass[10pt]{article}
\usepackage[a4paper,margin=19mm]{geometry}
\usepackage{amsmath,amssymb,mathtools}
\usepackage{booktabs,array,longtable}
\usepackage{enumitem}
\usepackage{xcolor}
\usepackage[hidelinks]{hyperref}

\newcommand{\Z}{\mathcal{Z}}
\newcommand{\e}{\mathrm{e}}
\newcommand{\unitstep}{\mathbf{1}}
\newcommand{\src}[1]{{\small\color{gray}[PDF pp.~#1]}}
\newcommand{\itemhead}[2]{\par\medskip\noindent\textbf{#1}\quad\src{#2}\par\nopagebreak}
\newcommand{\verify}[1]{\par\smallskip\noindent{\color{blue!55!black}\textbf{Check.}} #1\par}
\newcommand{\wrong}[1]{\par\smallskip\noindent{\color{red!55!black}\textbf{If you missed it.}} #1\par}
\setlist{nosep,leftmargin=*}
\setlength{\parskip}{0.35em}
\setlength{\parindent}{0pt}

\title{EE6203 Quiz 1 --- Evaluation Set\\\large Worked Solutions and Score Interpretation}
\author{Companion to \texttt{quizzes/quiz-01-evaluation-set.html}}
\date{31 August 2026}

\begin{document}
\maketitle

\section*{Before you open this}
\addcontentsline{toc}{section}{Before you open this}

This document contains the answers. It is separate from the question file for a
reason: the evaluation set is a \emph{measurement}, and a measurement taken with
the answer key in view measures nothing. Sit all 44 questions closed-book first,
with a calculator and Table 2--1 available, then come here.

Scope is Chapters 1--3 only, matching the Week~3 announcement that Quiz~1 is a
30-minute closed-book assessment excluding the Chapter~4 material. Every item is
new; none is reused from the drills, knowledge checks, worked archetypes or mock
quizzes in \texttt{notes/quiz1-prep-learning} or \texttt{notes/quiz1-prep-reviewing}.

Each numeric answer below was derived twice by independent routes --- a transform
argument and a direct time-domain or physical argument --- and the two agree. Where
the second route is quick enough to use under exam pressure, it is given as a
\textbf{Check}, because that is the habit worth carrying into the quiz.

\section*{How to read your topic scores}
\addcontentsline{toc}{section}{How to read your topic scores}

The per-topic bars matter more than the total. Treat a topic below 60\,\% as a gap,
60--79\,\% as shaky, and 80\,\%+ as sound. Then fix gaps in dependency order, not in
alphabetical order --- later topics are built out of earlier ones, and repairing
\textbf{H} while \textbf{C} is broken wastes the effort.

\begin{center}
\begin{tabular}{@{}cl>{\raggedright\arraybackslash}p{62mm}@{}}
\toprule
Tier & Topics & Why it comes first \\
\midrule
1 & B, C & The transform definition and the shift/translation theorems. Everything
        downstream is algebra on top of these. A gap here shows up as wrong answers
        in D, E, F and H that look like different mistakes but are not. \\
2 & D, E, F & Reading a transform back: legality of the final value, inverse methods,
        recurrences. These are the most heavily examined single-topic items. \\
3 & G, I & Impulse sampling, the starred transform, and the star-factorisation rule
        \(\left[GX^{*}\right]^{*}=G^{*}X^{*}\). This one rule governs every
        block-diagram question. \\
4 & H, J & ZOH-plus-plant conversion and the closed-loop configurations. Highest mark
        density in past questions, but only reachable once tiers 1--3 hold. \\
--- & A & Vocabulary. Cheap marks; revise by reading, not by drilling. \\
\bottomrule
\end{tabular}
\end{center}

\paragraph{Minimum viable path.} If time is short, the shortest route to a safe pass is
C2 (advance with initial conditions) \(\to\) D3/D4 (final-value legality) \(\to\)
E3 (repeated-pole inversion) \(\to\) H4 and H1 (the \((1-z^{-1})\Z\{G(s)/s\}\)
conversion) \(\to\) G2 (star factorisation) \(\to\) J1/J2 (which denominator).
Those seven skills, done reliably, cover the structure of every recalled past
question. Everything else in this set is either supporting arithmetic or vocabulary.

\paragraph{What you may safely defer.} A wrong answer on A2, B3 or G3 costs almost
nothing to repair and does not block anything else. Deferring E1 (method choice) is
also acceptable if D and E are otherwise strong, since the wrong method still
reaches the right answer, just slower.

\section*{Answer summary}
\addcontentsline{toc}{section}{Answer summary}

Mark yourself against this table first, then read only the items you missed.
Multiple-choice answers are quoted by content, because the question file shuffles
the option order.

\begin{longtable}{@{}llp{95mm}@{}}
\toprule
Item & Topic & Answer \\
\midrule
\endhead
A1 & A & Digital (discrete in time \emph{and} amplitude) \\
A2 & A & The discrete and continuous descriptions must be reconciled in one domain \\
A3 & A & \(26\) sample instants \\
\midrule
B1 & B & \(2+2z^{-1}+2z^{-2}-z^{-3}\) \\
B2 & B & \(z=\e^{-1}\approx 0.3679\) \\
B3 & B & \(T^{2}z^{-1}(1+z^{-1})/(1-z^{-1})^{3}\) \\
B4 & B & \(x[4]=0.5\) \\
\midrule
C1 & C & \(z^{-2}/(1-0.4z^{-1})\) \\
C2 & C & \(z^{2}X(z)-3z^{2}+z\) \\
C3 & C & Every \(z^{-1}\) scaled by \(\e^{-aT}\), so \(z^{-2}\) carries \(\e^{-2aT}\) \\
C4 & C & \(z=1.2\) \\
\midrule
D1 & D & \(x[0]=0\) \\
D2 & D & \(x[1]=2\) \\
D3 & D & Theorem applies; \(x[\infty]=0.5/1.9=0.263\) \\
D4 & D & Theorem does not apply; the remaining pole \(1.2\) is outside, so \(x[k]\) diverges \\
D5 & D & Two poles at \(z=1\) (double), so the theorem fails \\
\midrule
E1 & E & Direct division (long division) \\
E2 & E & \(x[2]=0.4375\) \\
E3 & E & \(x[k]=k(0.8)^{k-1}\) \\
E4 & E & \(x[3]=1.92\) \\
E5 & E & \(x[k]=\sin(0.5k)\) \\
\midrule
F1 & F & \(y[3]=2.56\) \\
F2 & F & \(Y(z)=z/[(z-0.5)(z-0.6)]\) \\
F3 & F & \(y[2]=1.1\) \\
F4 & F & \(y[k]=10\left[(0.6)^{k}-(0.5)^{k}\right]\) \\
\midrule
G1 & G & Impulse strength \(\e^{-1}\approx 0.3679\) \\
G2 & G & \(\left[G(s)X^{*}(s)\right]^{*}=G^{*}(s)X^{*}(s)\) \\
G3 & G & \(\omega_{s}=2\pi/T\approx 314.16\) rad/s \\
G4 & G & \(X(z)=1/(1-z^{-1})\) \\
\midrule
H1 & H & \(0.8647z^{-1}/(1-0.1353z^{-1})\) \\
H2 & H & \(G(1)=1\) \\
H3 & H & \(y[1]=1-\e^{-2}\approx 0.8647\) \\
H4 & H & \(X(z)=(1-z^{-1})\Z\{G(s)/s\}\) \\
H5 & H & \(g[1]\approx 0.2838\) \\
\midrule
I1 & I & \(0.1917\) \\
I2 & I & \(G(z)H(z)\neq\Z\{G(s)H(s)\}\); the intermediate sampler decides \\
I3 & I & No --- \(R\) cannot be separated from the dynamics \\
I4 & I & Causality of \(g(t)\) \\
\midrule
J1 & J & \(C(z)=G(z)R(z)/\left[1+GH(z)\right]\) \\
J2 & J & \(C(z)=G(z)R(z)/\left[1+G(z)H(z)\right]\) \\
J3 & J & \(G_{1}R(z)\) does not factor, so no input-independent ratio exists \\
J4 & J & \(K_{P}=1.8\) \\
J5 & J & \(G_{D}(z)=K_{P}+K_{I}/(1-z^{-1})+K_{D}(1-z^{-1})\) \\
J6 & J & \(m[0]=6.2\) \\
\bottomrule
\end{longtable}

\newpage
\section{Worked solutions}

\subsection{A --- Signals and the digital-control architecture}

\itemhead{A1.}{13--16}
A signal is \emph{sampled-data} when time is discrete but amplitude is still
continuous, and \emph{digital} when both are discrete. A 12-bit converter has
\(2^{12}\) representable levels, so the amplitude is quantised: the ADC output is
digital.
\wrong{You are treating ``sampled'' and ``digital'' as synonyms. The distinction is
the one Chapter~1 spends three slides on, and it is the reason quantisation is
ignored for the rest of the course --- with enough bits the digital signal is
\emph{modelled} as sampled-data, but that is an approximation, not an identity.}

\itemhead{A2.}{10, 16}
Figure 1.2 shows the loop crossing a boundary twice: the ADC turns the continuous
error \(e(t)\) into the sequence \(e[k]\), and the DAC-plus-ZOH turns the computed
\(u[k]\) back into a continuous \(u(t)\) that drives the plant. The controller is
therefore discrete while the plant remains continuous, and the two descriptions
cannot be multiplied together as they stand.

Reconciling them is precisely the work of the next two chapters: Chapter~2 builds the
discrete description (the \(z\) transform), and Chapter~3 converts the continuous
plant, together with its hold, into a pulse transfer function in the same variable.
Only then does a single algebra cover the whole loop.
\wrong{The clock synchronises the conversions; it does not abolish the boundary. And
nothing here asks for the plant to be redesigned --- the plant is whatever physics
gives you. If this framing is new, it is worth fixing early: every topology question
in I and J is an instance of it.}

\itemhead{A3.}{14--16}
\(0.1/0.004=25\) sampling intervals, hence \(k=0,1,\dots,25\), which is
\(\mathbf{26}\) sample instants.
\wrong{Answering 25 is the same off-by-one that corrupts piecewise-input transforms:
counting intervals instead of instants. On a half-open interval \(0\le t< nT\)
there are \(n\) samples; on a closed interval \(0\le t\le nT\) there are \(n+1\).}

\subsection{B --- Z-transform definition and elementary functions}

\itemhead{B1.}{24--25}
Tabulate before summing:
\[
\begin{array}{c|cccccc}
k & 0 & 1 & 2 & 3 & 4 & \ge 5\\\hline
x[k] & 2 & 2 & 2 & -1 & 0 & 0
\end{array}
\]
The interval \(0\le t<3T\) is half-open, so it holds \(k=0,1,2\); the sample at
\(k=3\) already belongs to the next piece. Substituting into
\(X(z)=\sum_{k\ge0}x[k]z^{-k}\),
\[
X(z)=2+2z^{-1}+2z^{-2}-z^{-3}.
\]
No table entry and no shift theorem is used; the coefficients \emph{are} the samples.

\itemhead{B2.}{32, entry 4}
Entry 4 reads \(\e^{-akT}\leftrightarrow 1/(1-\e^{-aT}z^{-1})\). With \(a=2\) and
\(T=0.5\), \(aT=1\), so the single pole is at \(z=\e^{-1}=0.3679\).
\verify{The sequence is \(x[k]=\e^{-2(0.5k)}=(\e^{-1})^{k}\), a geometric sequence
of ratio \(\e^{-1}\); a geometric sequence of ratio \(r\) has its pole at \(z=r\).}
\wrong{Answering \(-2\) or \(\e^{-2}\) confuses \(a\) with \(aT\). The pole depends
on the sampling period; the same continuous plant gives a different discrete pole
at a different \(T\).}

\itemhead{B3.}{32, entry 6}
\(\Z\{t^{2}\}\) sampled at period \(T\) is
\(T^{2}z^{-1}(1+z^{-1})/(1-z^{-1})^{3}\). The \((1+z^{-1})\) numerator factor is the
part most often lost. The entry \(Tz^{-1}/(1-z^{-1})^{2}\) is the ramp \(kT\).
\verify{Expand the first two coefficients: \(T^{2}z^{-1}(1+z^{-1})(1+3z^{-1}+\cdots)
= T^{2}z^{-1}+4T^{2}z^{-2}+\cdots\), matching \((kT)^{2}\) at \(k=1,2\):
\(T^{2}\) and \(4T^{2}\).}

\itemhead{B4.}{33, entry 20}
The pair is \(k\,a^{k-1}\leftrightarrow z^{-1}/(1-az^{-1})^{2}\). With \(a=0.5\) and
\(k=4\),
\[
x[4]=4\,(0.5)^{3}=0.5 .
\]
\wrong{Reading the pair as \(k a^{k}\) gives \(4(0.5)^{4}=0.25\). The \(k-1\) exponent
is what makes the sequence start at \(x[1]=1\) rather than \(x[1]=a\); checking the
first sample against the transform's leading coefficient catches this instantly.}

\subsection{C --- Properties and the translation theorems}

\itemhead{C1.}{35}
The shifting theorem gives \(\Z\{x[k-n]\unitstep[k-n]\}=z^{-n}X(z)\), so with
\(X(z)=1/(1-0.4z^{-1})\),
\[
\Z\{x[k-2]\unitstep[k-2]\}=\frac{z^{-2}}{1-0.4z^{-1}} .
\]
A delay changes \emph{where} the sequence starts; it does not move the pole.
\wrong{The variant \(0.16/(1-0.4z^{-1})\) applies the complex-translation theorem
instead, i.e.\ it scales the amplitude as \(a^{k}\) would. Keep the two apart:
\(z^{-n}\) shifts in time, \(X(z/a)\) scales in amplitude.}

\itemhead{C2.}{35--36}
For the one-sided transform,
\[
\Z\{x[k+n]\}=z^{n}\Bigl[X(z)-\sum_{j=0}^{n-1}x[j]z^{-j}\Bigr].
\]
With \(n=2\), \(x[0]=3\), \(x[1]=-1\):
\[
z^{2}\bigl[X(z)-3-(-1)z^{-1}\bigr]=z^{2}X(z)-3z^{2}+z .
\]
\wrong{The sign of the last term is the whole question. An advance \emph{removes}
the samples the sequence has already produced, so each removed sample enters with a
minus sign --- and a negative sample therefore enters with a plus. Every forced
recurrence in topic F depends on getting this right.}

\itemhead{C3.}{39--40; 33, entry 16}
The complex-translation theorem states \(\Z\{\e^{-akT}x[k]\}=X(\e^{aT}z)\), which in
terms of \(z^{-1}\) means: replace every \(z^{-1}\) by \(\e^{-aT}z^{-1}\). Applying
that to the sine pair,
\[
\frac{\e^{-aT}z^{-1}\sin\omega T}
     {1-2\e^{-aT}z^{-1}\cos\omega T+\e^{-2aT}z^{-2}},
\]
where the \(z^{-2}\) term picks up \(\e^{-2aT}\) because it carries two factors of
\(z^{-1}\). This is entry 16 of Table 2--1.
\wrong{Scaling only the numerator is a half-application. The rule is mechanical ---
substitute, do not patch --- and the \(\e^{-2aT}\) on the \(z^{-2}\) term is the
marker that it was done properly.}

\itemhead{C4.}{34, 39}
\(\Z\{a^{k}x[k]\}=X(z/a)\). A pole of the new transform occurs where the argument
equals an old pole: \(z/2=0.6\), hence \(z=1.2\).
\wrong{Dividing instead of multiplying gives 0.3. Note what happened physically:
multiplying a sequence by \(2^{k}\) pushed a stable pole outside the unit circle.
That is the same mechanism that makes \(|a|\ge1\) forcing terms diverge in topic D.}

\subsection{D --- Initial- and final-value theorems}

\itemhead{D1--D2.}{42, 51--53}
The initial-value theorem is \(x[0]=\lim_{z\to\infty}X(z)\). Every numerator term of
\[
X(z)=\frac{2z^{-1}+3z^{-2}}{1-0.5z^{-1}+0.06z^{-2}}
\]
carries a negative power of \(z\), so the numerator vanishes as \(z\to\infty\) while
the denominator tends to 1. Hence \(x[0]=0\).

For \(x[1]\), multiply out and match coefficients of \(z^{-1}\):
\[
x[1]-0.5\,x[0]=2 \quad\Longrightarrow\quad x[1]=2 .
\]
\verify{Continue the recursion, matching each coefficient in turn:
\[
x[k]=\text{(numerator)}_{k}+0.5\,x[k-1]-0.06\,x[k-2]
\;\Longrightarrow\;
0,\;2,\;4,\;1.88,\;0.70,\;0.2372,\dots
\]
The sequence decays, as the poles at \(0.2\) and \(0.3\) require.}
\wrong{Reading the leading numerator coefficient as \(x[0]\) gives 2, which is
actually \(x[1]\). The initial-value theorem is about the \(z^{0}\) coefficient, and
this numerator has none.}

\itemhead{D3.}{44--46}
The theorem (p.~44) requires all poles of \(X(z)\) inside the unit circle,
\emph{with the possible exception of a simple pole at \(z=1\)}. Form
\[
(1-z^{-1})X(z)=\frac{0.5z^{-1}}{1+0.9z^{-1}} ,
\]
whose only pole is \(z=-0.9\). Since \(|-0.9|<1\), the theorem applies and
\[
x[\infty]=\lim_{z\to1}\frac{0.5z^{-1}}{1+0.9z^{-1}}=\frac{0.5}{1.9}=0.263 .
\]
\wrong{Rejecting the theorem because \(-0.9\) is negative confuses sign with
magnitude. A negative pole gives an \emph{alternating} transient, and an alternating
transient still decays provided its magnitude is below one. Only \(|{\cdot}|=1\)
(sustained alternation, as at \(z=-1\)) or \(>1\) destroys the limit.}

\itemhead{D4.}{44--46}
Removing the step pole,
\[
(1-z^{-1})X(z)=\frac{z-1}{z}\cdot\frac{z}{(z-1)(z-1.2)}=\frac{1}{z-1.2},
\]
which has its pole at \(z=1.2\), outside the unit circle. The theorem therefore does
not apply, and indeed \(x[k]\) contains a \((1.2)^{k}\) mode and diverges. The value
\(-5\) obtained by substitution is arithmetically correct and physically meaningless.
\wrong{This is the pitfall the Remarks slide (p.~46) warns about in exactly these
words: the theorem ``can yield misleading results'' for unbounded or oscillatory
sequences. The discipline is fixed --- cancel the permitted pole at \(z=1\) first,
classify what survives, and only then take the limit. Never take the limit first.}

\itemhead{D5.}{32 entry 5; 44}
\(X(z)=Tz^{-1}/(1-z^{-1})^{2}\) has a \emph{double} pole at \(z=1\). The theorem
permits at most a simple one, so it does not apply --- correctly, since \(x[k]=kT\)
grows without bound.
\wrong{If you answered 1, you read \((1-z^{-1})^{2}\) as a single factor. The
exponent is the pole multiplicity, and multiplicity is exactly what separates the
step (allowed) from the ramp (not allowed).}

\subsection{E --- Inverse z-transform methods}

\itemhead{E1.}{49--53}
Long division produces \(x[0],x[1],x[2],\dots\) directly, one coefficient per step,
and never yields a general term. Partial-fraction expansion of \(X(z)/z\) yields the
closed form but costs a factorisation and a residue calculation. With only four
samples wanted and no closed form required, long division is faster.
\wrong{Choosing partial fractions here is not \emph{wrong}, only slow --- and in a
30-minute closed-book quiz, slow is a mark. Match the method to what is asked for:
``find \(x[k]\)'' means PFE, ``find the first few samples'' means long division.}

\itemhead{E2.}{60--62}
Write \(u=z^{-1}\) and expand
\(1/[(1-0.5u)(1-0.25u)]=A/(1-0.5u)+B/(1-0.25u)\).
Evaluating at \(u=2\) gives \(A=2\); at \(u=4\) gives \(B=-1\). Hence
\(x[k]=2(0.5)^{k}-(0.25)^{k}\), and
\[
x[2]=2(0.25)-0.0625=0.4375 .
\]
\verify{Long division of \(1/(1-0.75z^{-1}+0.125z^{-2})\) gives \(1,\;0.75,\;0.4375,\;
0.234375\), matching the closed form term by term. Two independent routes agreeing is
the standard for anything you will write down under time pressure.}

\itemhead{E3--E4.}{33 entry 20; 61}
Convert to the table's variable before matching:
\[
X(z)=\frac{z}{(z-0.8)^{2}}
=\frac{z^{-1}}{(1-0.8z^{-1})^{2}}
\;\longleftrightarrow\; x[k]=k\,(0.8)^{k-1}.
\]
Then \(x[3]=3(0.8)^{2}=1.92\).
\verify{Long division of \(z^{-1}/(1-1.6z^{-1}+0.64z^{-2})\) gives \(0,\;1,\;1.6,\;
1.92,\;2.048\). Note the sequence \emph{rises} before decaying --- the \(k\) factor
beats the \(0.8^{k}\) decay until \(k\approx4.5\), which is what a repeated pole
inside the unit circle looks like.}
\wrong{The two index-shifted variants \(k(0.8)^{k}\) and \((k+1)(0.8)^{k}\) both give
the wrong \(x[1]\): the true value is 1, they give 0.8 and 1.6. Testing a candidate
closed form at \(k=0\) and \(k=1\) costs five seconds and settles the question.}

\itemhead{E5.}{33 entry 14; 63--64}
Entry 14 is
\(\sin\omega kT \leftrightarrow z^{-1}\sin\omega T/(1-2z^{-1}\cos\omega T+z^{-2})\).
Comparing, \(\omega T=0.5\), so \(x[k]=\sin(0.5k)\).
\wrong{A cosine would show \(1-z^{-1}\cos\omega T\) in the numerator; a damped sine
would carry \(\e^{-aT}\) factors (entry 16, as in C3). Recognising the shape of the
denominator \(1-2z^{-1}\cos\omega T+z^{-2}\) --- unit leading and trailing
coefficients, poles on the unit circle --- is faster than any expansion.}

\subsection{F --- Difference equations solved by z transform}

\itemhead{F1.}{67--69}
Transform the homogeneous equation with its initial condition:
\[
zY(z)-zy[0]-0.8Y(z)=0
\;\Longrightarrow\;
Y(z)=\frac{5z}{z-0.8},
\qquad y[k]=5(0.8)^{k}.
\]
Hence \(y[3]=5(0.512)=2.56\).
\verify{Iterate the rule directly: \(5,\;4,\;3.2,\;2.56\).}
\wrong{If a \(1/(1-z^{-1})\) factor appeared in your working, you added a forcing
term that is not there. An unforced recurrence has no step pole and therefore no
constant steady-state component; it decays to zero.}

\itemhead{F2--F4.}{35--36, 60, 67--69}
For a two-step advance,
\(\Z\{y[k+2]\}=z^{2}Y(z)-z^{2}y[0]-zy[1]\). With \(y[0]=0\), \(y[1]=1\):
\[
z^{2}Y-z-1.1zY+0.3Y=0
\;\Longrightarrow\;
Y(z)=\frac{z}{z^{2}-1.1z+0.3}=\frac{z}{(z-0.5)(z-0.6)} .
\]
The factorisation checks out because \(0.5+0.6=1.1\) and \(0.5\times0.6=0.3\).

Directly from the rule, \(y[2]=1.1y[1]-0.3y[0]=1.1\).

For the closed form, expand \(Y(z)/z\):
\[
\frac{Y(z)}{z}=\frac{1}{(z-0.5)(z-0.6)}
=\frac{-10}{z-0.5}+\frac{10}{z-0.6},
\qquad
y[k]=10\left[(0.6)^{k}-(0.5)^{k}\right].
\]
\verify{Test the closed form at three points: \(y[0]=10(1-1)=0\) \checkmark,
\(y[1]=10(0.1)=1\) \checkmark, \(y[2]=10(0.36-0.25)=1.1\) \checkmark --- matching the
initial conditions and the sample generated from the recurrence. The full opening is
\(0,\;1,\;1.1,\;0.91,\;0.671,\;0.4651\), decaying as both poles are inside the unit
circle.}
\wrong{A numerator of 1 or \(z^{2}\) in F2 means the initial-condition terms were
mishandled --- go back to C2, which is the same rule in isolation. In F4, the
three-point test rejects every wrong option in under a minute, so there is no excuse
for guessing on a closed form.}

\subsection{G --- Impulse sampling and the starred transform}

\itemhead{G1.}{75--79}
The impulse-sampled signal is
\(x^{*}(t)=\sum_{k\ge0}x(kT)\delta(t-kT)\). With \(T=0.5\), the impulse at
\(t=1.0\) is the one with \(k=2\), and its \emph{strength} (area) is
\(x(1.0)=\e^{-1}=0.3679\).
\wrong{An ideal impulse has zero width and infinite height, so there is no amplitude
to read; the coefficient multiplying \(\delta(t-kT)\) is an area. This is not
pedantry --- it is why the ZOH must be modelled separately as
\((1-\e^{-Ts})/s\) rather than being absorbed into the sampler.}

\itemhead{G2.}{103--106}
The one general identity is
\[
\bigl[G(s)X^{*}(s)\bigr]^{*}=G^{*}(s)X^{*}(s).
\]
A factor that is \emph{already} an impulse train can be pulled out of the star,
because sampling an already-sampled signal changes nothing. An unsampled product
\(G(s)H(s)\) or \(G(s)X(s)\) cannot be split: there is no sampler between the
factors to break the convolution.
\wrong{This single rule decides topics I and J almost entirely. If it is not
automatic, do nothing else until it is: it is why J1 has \(GH(z)\) in the denominator
while J2 has \(G(z)H(z)\), and why I3 has no pulse transfer function at all.}

\itemhead{G3.}{77--78, 103--104}
\(X^{*}(s)\) is periodic along the imaginary axis with
\(\omega_{s}=2\pi/T\). For \(T=0.02\) s, \(\omega_{s}=2\pi/0.02=314.16\) rad/s.
\wrong{Answering 50 gives \(f_{s}=1/T\) in hertz. The factor \(2\pi\) is the whole
question, and the same factor reappears wherever \(\omega\) and \(f\) are exchanged.}

\itemhead{G4.}{78, 103--104}
Impulse-sampling a unit step gives
\(X^{*}(s)=\sum_{k\ge0}\e^{-kTs}=1/(1-\e^{-Ts})\) as a geometric series.
The definition \(z=\e^{Ts}\) turns \(\e^{-Ts}\) into \(z^{-1}\), so
\[
X(z)=\frac{1}{1-z^{-1}} ,
\]
which is Table 2--1 entry 3, as it must be.
\wrong{This is the bridge between the two descriptions: \(X^{*}(s)\) and \(X(z)\)
are the same object in different variables, which is why a single sampled block may
be written either way. It is also why \(z=\e^{Ts}\) may \emph{not} be substituted
into an ordinary continuous plant \(G_{p}(s)\) --- the substitution applies to
starred quantities, not to unsampled ones.}

\subsection{H --- Data holds and ZOH-plus-plant discretisation}

\itemhead{H1--H3.}{111--115}
Equation (3--49) gives the ZOH-plus-plant conversion. With
\(G_{p}(s)=4/(s+4)\) and \(T=0.5\),
\begin{align*}
G(z)&=(1-z^{-1})\,\Z\!\left\{\frac{4}{s(s+4)}\right\}
     =(1-z^{-1})\,\Z\!\left\{\frac{1}{s}-\frac{1}{s+4}\right\}\\
    &=(1-z^{-1})\left[\frac{1}{1-z^{-1}}-\frac{1}{1-\e^{-2}z^{-1}}\right]
     =1-\frac{1-z^{-1}}{1-\e^{-2}z^{-1}}
     =\frac{(1-\e^{-2})z^{-1}}{1-\e^{-2}z^{-1}} .
\end{align*}
With \(\e^{-2}=0.1353\), \(G(z)=0.8647z^{-1}/(1-0.1353z^{-1})\).

The DC gain is \(G(1)=0.8647/(1-0.1353)=1\), matching \(G_{p}(0)=1\).

For a sampled unit step, \(Y(z)=G(z)/(1-z^{-1})\), whose \(z^{-1}\) coefficient is
\(0.8647\); so \(y[0]=0\) and \(y[1]=0.8647\).
\verify{The ZOH holds the command constant over one period, so \(y(T)\) is just the
continuous step response of the plant at \(t=T\): \(1-\e^{-4(0.5)}=1-\e^{-2}=0.8647\).
The transform and the physics agree, and the agreement costs one line to check.}
\wrong{Losing the \(z^{-1}\) is the most common error and the most consequential:
a ZOH driving a strictly proper plant \emph{always} produces a one-sample delay,
because the output cannot respond before the first hold interval has elapsed. If
your \(G(z)\) has a nonzero \(g[0]\), something is wrong before you go any further.}

\itemhead{H4.}{113--114}
Equation (3--49) reads
\[
X(s)=\frac{1-\e^{-Ts}}{s}G(s)
\quad\Longrightarrow\quad
X(z)=(1-z^{-1})\,\Z\!\left\{\frac{G(s)}{s}\right\}.
\]
Only the numerator factor \(1-\e^{-Ts}\) becomes \(1-z^{-1}\) (it is a one-sample
delay, and \(\e^{-Ts}\leftrightarrow z^{-1}\)). The \(1/s\) belongs to the hold's own
integration and must stay \emph{inside} the transform alongside \(G(s)\).
\wrong{The variant \((1-z^{-1})\Z\{G(s)\}\) drops that \(1/s\) and silently
discretises the wrong plant. The variant keeping \(1-\e^{-Ts}\) outside leaves an
\(s\)-domain factor multiplying a \(z\)-domain expression, which is not a
well-formed object at all --- a useful sanity filter in itself.}

\itemhead{H5.}{113--115}
With \(G_{p}(s)=1/[s(s+2)]\) and \(T=1\),
\[
\frac{G_{p}(s)}{s}=\frac{1}{s^{2}(s+2)}
=-\frac{0.25}{s}+\frac{0.5}{s^{2}}+\frac{0.25}{s+2},
\]
(coefficients from \(1=As(s+2)+B(s+2)+Cs^{2}\) at \(s=0\), \(s=-2\), and matching
\(s^{2}\)). Therefore
\[
G(z)=(1-z^{-1})\left[\frac{-0.25}{1-z^{-1}}
+\frac{0.5z^{-1}}{(1-z^{-1})^{2}}
+\frac{0.25}{1-\e^{-2}z^{-1}}\right].
\]
Collecting the \(z^{-1}\) coefficient: the first term contributes 0, the second
\(0.5\), and the third \(0.25(\e^{-2}-1)=-0.2162\). Hence \(g[1]=0.2838\).
\verify{Independently, \(g[k]=s_{r}(kT)-s_{r}((k-1)T)\) where
\(s_{r}(t)=-0.25+0.5t+0.25\e^{-2t}\) is the plant's continuous step response. Then
\(s_{r}(0)=0\) and \(s_{r}(1)=0.2838\), so \(g[0]=0\) and \(g[1]=0.2838\). The two
routes agree, and \(g[0]=0\) confirms the expected one-sample delay.}

\subsection{I --- Pulse transfer functions, samplers and cascades}

\itemhead{I1--I2.}{117--121}
Without an intermediate sampler the blocks must be combined \emph{before} sampling:
\[
GH(s)=\frac{1}{(s+1)(s+3)}
\;\Longrightarrow\;
gh(t)=\tfrac12\left(\e^{-t}-\e^{-3t}\right).
\]
Sampling at \(T=0.5\) gives \(gh(0)=0\) and
\(gh(0.5)=\tfrac12(0.6065-0.2231)=0.1917\).

With an intermediate sampler, each block gets its own pulse transfer function and
the sequences convolve: \(g[k]=\e^{-0.5k}\), \(h[k]=\e^{-1.5k}\), so the product's
coefficients are \(1,\;0.8297,\;0.5530,\dots\)

The two disagree at \emph{every} coefficient, most starkly at \(k=0\), where one is 0
and the other is 1. They are different systems, and no choice of \(T\) reconciles
them.
\wrong{The \(k=0\) comparison is the fastest possible proof and worth remembering as
a stock argument: \(GH(z)\) is strictly proper (two poles, no zeros, so \(g h[0]=0\)),
while \(G(z)H(z)\) is a product of two transforms each starting at 1. If a
block-diagram answer ever hinges on which form applies, this settles it in one line.}

\itemhead{I3.}{107--110, 126}
The signal reaching the first sampler is \(F(s)R(s)\), so the sampled quantity is
\(\left[F(s)R(s)\right]^{*}\). By G2, a star may be factored only when one factor is
already starred, and neither is here. Consequently \(R\) cannot be separated from the
dynamics and no input-independent ratio \(C(z)/R(z)\) exists.

Contrast the repaired configuration: with a sampler placed \emph{before} \(F(s)\),
\(\left[F(s)X^{*}(s)\right]^{*}=F^{*}(s)X^{*}(s)\), and \(Y(z)=F(z)X(z)\) does exist.
\wrong{Sampling the output does not repair a missing input sampler --- it only says
where the continuous output is observed. Also note what is \emph{not} being claimed:
the response to a specified \(R(s)\) remains perfectly computable. Only the
input-independent ratio fails to exist.}

\itemhead{I4.}{95--100}
The plant impulse response satisfies \(g(t)=0\) for \(t<0\), so
\(g\bigl((k-h)T\bigr)=0\) whenever \(h>k\) and those terms drop out of the sum. This
is causality, and nothing else: an unstable but causal plant has the same summation
limits, and a hold changes what \(g\) is, not where the sum stops.

\subsection{J --- Closed-loop configurations, digital controller and PID}

\itemhead{J1.}{122--125}
Case 1 of Table 3--1. Derive it rather than recalling it:
\[
E(s)=R(s)-H(s)C(s),\qquad C(s)=G(s)E^{*}(s).
\]
Substituting, \(E(s)=R(s)-H(s)G(s)E^{*}(s)\). Starring both sides and using
\(\left[HGE^{*}\right]^{*}=\left[HG\right]^{*}E^{*}\):
\[
E^{*}=R^{*}-\left[GH\right]^{*}E^{*}
\;\Longrightarrow\;
E(z)=\frac{R(z)}{1+GH(z)},
\qquad
C(z)=\frac{G(z)R(z)}{1+GH(z)} .
\]
There is no sampler between \(G\) and \(H\), so the loop term stays welded together
as \(GH(z)=\Z\{G(s)H(s)\}\). The forward path \emph{is} driven by \(E^{*}\), which is
why the numerator does separate.

\itemhead{J2.}{124--125}
Case 2. Now \(H\) is driven by the sampled output, so \(E=R-HC^{*}\) and
\(\left[HC^{*}\right]^{*}=H^{*}C^{*}\):
\[
E^{*}=R^{*}-H^{*}G^{*}E^{*}
\;\Longrightarrow\;
C(z)=\frac{G(z)R(z)}{1+G(z)H(z)} .
\]
\wrong{J1 and J2 have identical numerators and differ only in whether the denominator
reads \(GH(z)\) or \(G(z)H(z)\). One sampler decides it. Before writing any
closed-loop formula, circle every sampler on the diagram and ask, for each product,
whether a sampler separates the two blocks; if not, they stay inside one \(\Z\{\cdot\}\).}

\itemhead{J3.}{125--126}
Case 4 has no sampler after the summing junction, so \(G_{1}\) is driven by the
continuous error. Writing \(e=R-HC\) and sampling the signal that actually reaches
the sampler, \(\left[G_{1}e\right]^{*}\), gives
\[
\left[G_{1}e\right]^{*}=\left[G_{1}R\right]^{*}-\left[G_{1}G_{2}H\right]^{*}\left[G_{1}e\right]^{*},
\qquad
C(z)=\frac{G_{2}(z)\,G_{1}R(z)}{1+G_{1}G_{2}H(z)} .
\]
The term \(G_{1}R(z)=\Z\{G_{1}(s)R(s)\}\) does not factor into \(G_{1}(z)R(z)\),
because \(R\) passes through \(G_{1}\) before any sampler. A pulse transfer function
is by definition an input-independent ratio, so this system has none --- the point
the lecture makes explicitly on p.~126.
\wrong{Note precisely what fails. The system is not unstable, not unanalysable, and
not lacking an output: for any \emph{specified} \(R(s)\) the response is computable
by the same technique. What does not exist is the ratio.}

\itemhead{J4, J6.}{134--135}
With \(K=2\), \(T=0.1\), \(T_{i}=0.5\), \(T_{d}=0.2\):
\[
K_{I}=\frac{KT}{T_{i}}=0.4,\qquad
K_{D}=\frac{KT_{d}}{T}=4,\qquad
K_{P}=K-\frac{KT}{2T_{i}}=K-\frac{K_{I}}{2}=2-0.2=1.8 .
\]
For a unit discrete impulse at the error input, read the \(k=0\) coefficient off each
term of \(G_{D}(z)\):
\[
m[0]=K_{P}+K_{I}+K_{D}=1.8+0.4+4=6.2,
\qquad
m[1]=K_{I}-K_{D}=-3.6 .
\]
\wrong{Answering \(K_{P}=2\) means the \(-K_{I}/2\) correction was missed. It is not
cosmetic: the trapezoidal rule averages \(e[k-1]\) and \(e[k]\), and half of the
current sample's contribution is already counted inside the integral term, so the
proportional gain must be reduced to compensate. The slide states \(K_{P}<K\)
explicitly, which is a free check on your arithmetic. The large \(|m[1]|\) is the
derivative kick, and it is why \(K_{D}\) is the term that makes PID noise-sensitive.}

\itemhead{J5.}{135}
\[
G_{D}(z)=\frac{M(z)}{E(z)}=K_{P}+\frac{K_{I}}{1-z^{-1}}+K_{D}\left(1-z^{-1}\right).
\]
Integration accumulates, giving a pole at \(z=1\); differentiation differences,
giving a zero at \(z=1\). Swapping the two --- \(K_{I}(1-z^{-1})\) with
\(K_{D}/(1-z^{-1})\) --- inverts both roles at once.
\wrong{A form with \(T\) still visible has double-counted the sampling period:
\(T\) is already absorbed into \(K_{I}=KT/T_{i}\) and \(K_{D}=KT_{d}/T\). Reading the
gain definitions and the transfer function as one package prevents this.}

\section*{After marking}
\addcontentsline{toc}{section}{After marking}

Write your ten topic percentages down before doing anything else, then act on them
in the tier order given at the front. Two habits are worth carrying into the quiz
regardless of the scores, because they appear in the \textbf{Check} lines above more
often than any single formula:

\begin{enumerate}[label=\arabic*.]
\item \textbf{Generate a sample.} Any closed form, inverse transform or recurrence
solution can be tested at \(k=0\) and \(k=1\) against the original time-domain rule
in well under a minute. Every wrong option in B4, E3, F4 and D2 dies to this test.
\item \textbf{Check a gain or a limit physically.} A ZOH-plus-plant conversion must
reproduce the plant's DC gain; a stable sampled loop must settle where the continuous
one would. H2 and H3 are that check applied deliberately, and it catches a lost
\((1-z^{-1})\) factor faster than re-deriving the algebra.
\end{enumerate}

Where a score is low and you are unsure whether the cause is the concept or the
arithmetic, re-attempt only the items you missed --- the question file keeps the
option order fixed, so a second pass measures the same thing.

\end{document}
